% Part: first-order-logic
% Chapter: models-theories
% Section: set-theory

\documentclass[../../../include/open-logic-section]{subfiles}

\begin{document}

\olfileid{fol}{mat}{set}

\olsection{نظرية المجموعات}

يمكن تطوير الرياضيات كلها تقريبًا في نظرية المجموعات. وينطوي تطوير
الرياضيات في هذه النظرية على أمور عدة. فأولًا، يتطلب مجموعة من
البديهيات للعلاقة~$\in$. وقد طُورت أنساق بديهيات مختلفة، تكون
خصائص~$\in$ فيها متعارضة أحيانًا. ويبرز بينها نسق البديهيات المعروف
بـ~$\Log{ZFC}$، أي نظرية المجموعات لزرميلو--فرينكل مع بديهية
الاختيار: فهو أكثرها استعمالًا ودراسة بفارق كبير، إذ يتبين أن
بديهياته تكفي لإثبات كل ما يتوقع الرياضيون تقريبًا إمكان إثباته. لكن
قبل إقامة ذلك، لا بد أولًا من توضيح كيف نستطيع أصلًا أن
\emph{نعبر} عن كل ما يود الرياضيون التعبير عنه. فاللغة، بادئ ذي بدء،
لا تحتوي !!{constant}s ولا !!{function}s، ولذلك لا يبدو واضحًا لأول
وهلة أننا نستطيع الحديث عن مجموعات بعينها (مثل $\emptyset$ أو
$\Nat$)، أو عن عمليات على المجموعات (مثل $X \cup Y$ و$\Pow{X}$)،
فضلًا عن أبنية أخرى تدخل فيها أشياء غير المجموعات، كالعلاقات والدوال.

وليست علاقة «هو عنصر من» العلاقة الوحيدة التي تعنينا؛ فعلاقة «هو
مجموعة جزئية من» لا تكاد تقل عنها أهمية. ولكن يمكننا \emph{تعريف}
«هو مجموعة جزئية من» بواسطة «هو عنصر من». ولفعل ذلك، علينا أن نجد
!!a{formula}~$!A(x, y)$ في لغة نظرية المجموعات يرضيها زوج المجموعتين
~$\tuple{X, Y}$ إذا وفقط إذا كان $X \subseteq Y$. ولكن $X$ مجموعة
جزئية من $Y$ بالضبط حين تكون جميع !!{element}s~$X$ !!{element}s في
~$Y$ أيضًا. ولذلك يمكننا تعريف $\subseteq$ بالصيغة
\[
\lforall[z][(z \in x \lif z \in y)]
\]
والآن، كلما أردنا استعمال العلاقة~$\subseteq$ في صيغة، استطعنا أن
نستعمل بدلًا منها تلك الصيغة (بعد إحلال المناسب محل $x$ و$y$، وإعادة
تسمية المتغير المقيد~$z$ إن لزم الأمر). فامتدادية المجموعات تعني،
مثلًا، أنه إذا كانت أي مجموعتين~$x$ و$y$ محتواةً كل منهما في الأخرى،
وجب أن تكونا المجموعة نفسها. ويمكن التعبير عن ذلك بـ
$\lforall[x][\lforall[y][((x \subseteq y \land y
    \subseteq x) \lif x = y)]]$، أو، إذا استبدلنا $\subseteq$
بالتعريف أعلاه، بـ
\[
\lforall[x][\lforall[y][((\lforall[z][(z \in x \lif z \in y)] \land
    \lforall[z][(z \in y \lif z \in x)]) \lif x = y)]].
\]
وهذه في الواقع إحدى بديهيات~$\Log{ZFC}$، وهي «بديهية الامتدادية».

لا يوجد !!{constant} لـ~$\emptyset$، لكننا نستطيع التعبير عن «$x$
خالية» بـ~$\lnot \lexists[y][y \in x]$. وعندئذ تصبح عبارة
«$\emptyset$ موجودة» ال!!{sentence}~$\lexists[x][\lnot
\lexists[y][y \in x]]$. وهذه بديهية أخرى من بديهيات~$\Log{ZFC}$.
(لاحظ أن بديهية الامتدادية تقتضي وجود مجموعة خالية واحدة فقط.) وكلما
أردنا الحديث عن $\emptyset$ في لغة نظرية المجموعات، كتبنا ذلك على
صورة «توجد مجموعة خالية و\dots». وللتعبير، مثلًا، عن أن $\emptyset$
مجموعة جزئية من كل مجموعة، يمكننا أن نكتب
\[
\lexists[x][(\lnot\lexists[y][y \in x] \land \lforall[z][x \subseteq
    z])]
\]
على أن $x \subseteq z$ يجب، بالطبع، أن يستبدل بدوره بتعريفه.

وللحديث عن عمليات على المجموعات، مثل $X \cup Y$ و$\Pow{X}$، علينا
أن نستعمل حيلة مماثلة. فلا رموز دوال في لغة نظرية المجموعات، لكننا
نستطيع التعبير عن العلاقتين الداليتين $X \cup Y = Z$ و$\Pow{X} = Y$
بالصيغتين
\begin{align*}
& \lforall[u][((u \in x \lor u \in y) \liff u \in z)]\\
& \lforall[u][(u \subseteq x \liff u \in y )]
\end{align*}
لأن !!{element}s $X \cup Y$ هي بالضبط المجموعات التي تكون
!!{element}s في~$X$ أو !!{element}s في~$Y$، و!!{element}s
$\Pow{X}$ هي بالضبط المجموعات الجزئية من~$X$. غير أن ذلك لا يتيح لنا
استعمال $x \cup y$ أو $\Pow{x}$ كما لو كانا حدين؛ إذ لا يمكننا إلا
استعمال ال!!{formula}s الكاملة التي تعرّف العلاقتين $X \cup Y = Z$
و$\Pow{X} = Y$. والواقع أننا لا نعرف أن هاتين العلاقتين تُرضيان في
أي حال، أي لا نعرف أن الاتحادات ومجموعات القوى موجودة دائمًا. فمثلًا،
ال!!{sentence} $\lforall[x][\lexists[y][\Pow{x} = y]]$ بديهية أخرى
من بديهيات~$\Log{ZFC}$ (بديهية مجموعة القوى).

وماذا عن الحديث عن الأزواج المرتبة أو الدوال؟ علينا هنا أن نشرح كيف
يمكننا التفكير في الأزواج المرتبة والدوال بوصفها أنواعًا خاصة من
المجموعات. ومن سبل تعريف الزوج المرتب $\tuple{x, y}$ أن نعرّفه بأنه
المجموعة $\{\{x\}, \{x, y\}\}$. ولكننا، كما سبق، لا نستطيع إدخال
!!a{function} تسمي هذه المجموعة؛ وإنما نستطيع تعريف العلاقة
$\tuple{x, y} = z$، أي $\{\{x\}, \{x, y\}\} = z$:
\[
\lforall[u][(u \in z \liff (\lforall[v][(v \in u \liff v = x)] \lor
  \lforall[v][(v \in u \liff (v = x \lor v = y))]))]
\]
وتقول هذه الصيغة إن !!{element}s~$u$ في~$z$ هي بالضبط المجموعات التي
يكون !!{element}ها الوحيد هو $x$، أو تكون !!{element}sها الوحيدة هي
$x$ و~$y$ (وبعبارة
أخرى، المجموعات المطابقة إما لـ~$\{x\}$ وإما لـ~$\{x, y\}$). وبعد
ذلك نستطيع قول أشياء أخرى، منها أن $X \times Y = Z$:
\[
\lforall[z][(z \in Z \liff \lexists[x][\lexists[y][(x \in
        X \land y \in Y \land \tuple{x, y} = z)]])]
\]

يمكن التفكير في الدالة $f \colon X \to Y$ بوصفها العلاقة $f(x)
= y$، أي بوصفها مجموعة الأزواج~$\Setabs{\tuple{x,y}}{f(x) = y}$.
وعندئذ نستطيع القول إن مجموعة~$f$ دالة من $X$ إلى $Y$ إذا كانت (أ)
علاقة $\subseteq X \times Y$، و(ب) كلية، أي إنه لكل $x \in X$ يوجد
$y \in Y$ ما بحيث $\tuple{x, y} \in f$، و(ج) دالية، أي إنه متى كان
$\tuple{x, y}, \tuple{x, y'} \in f$، كان $y = y'$ (لأن قيم الدوال
يجب أن تكون وحيدة). ولذلك يمكن كتابة «$f$ دالة من $X$ إلى $Y$» هكذا:
\begin{align*}
 \lforall[u][(u \in f \lif {}] & \lexists[x][\lexists[y][(x \in X \land y \in
      Y \land \tuple{x, y} = u)]]) \land {}\\
 \lforall[x][(x \in X \lif {}] &
      (\lexists[y][(y \in Y \land \mathrm{maps}(f, x, y))] \land {}\\
& (\lforall[y][\lforall[y'][((\mathrm{maps}(f, x, y) \land
    \mathrm{maps}(f, x, y')) \lif y = y')]]))
\end{align*}
حيث تختصر $\mathrm{maps}(f, x, y)$ الصيغة $\lexists[v][(v \in f \land
  \tuple{x, y} = v)]$ (وهذه ال!!{formula} تعبر عن «$f(x) = y$»).

وليس من الصعب الآن، مثلًا، التعبير عن أن $f\colon X \to Y$
!!{injective}:
\begin{multline*}
 f \colon X \to Y \land \lforall[x][\lforall[x'][((x \in X \land x' \in
     X \land {}]] \\
 \lexists[y][(\mathrm{maps}(f, x, y) \land \mathrm{maps}(f,
        x', y))]) \lif x = x')
\end{multline*}
تكون الدالة~$f\colon X \to Y$ !!{injective} إذا وفقط إذا كان
$x = x'$ كلما أرسلت $x, x' \in X$ إلى قيمة واحدة~$y$. وإذا اختصرنا
هذه الصيغة إلى $\mathrm{inj}(f, X, Y)$، كنا قد بلغنا موضعًا نستطيع فيه
صياغة قضية غير هينة، مثل مبرهنة كانتور، بلغة نظرية المجموعات: لا توجد
دالة !!{injective} من $\Pow{X}$ إلى $X$:
\[
\lforall[X][\lforall[Y][(\Pow{X} = Y \lif
    \lnot\lexists[f][\mathrm{inj}(f, Y, X)])]]
\]

قد يظن المرء أن نظرية المجموعات تحتاج إلى بديهية أخرى تضمن وجود
مجموعة لكل خاصية معرِّفة. فإذا كانت $!A(x)$ صيغة في نظرية المجموعات
يقع فيها المتغير~$x$ حرًا، أمكننا أن ننظر في ال!!{sentence}
\[
\lexists[y][\lforall[x][(x \in y \liff !A(x))]].
\]
وتقرر هذه ال!!{sentence} أن ثمة مجموعة~$y$ تكون !!{element}sها كل
قيم $x$ التي ترضي~$!A(x)$ وتلك القيم وحدها. ويسمى هذا المخطط «مبدأ
الاستيعاب». وهو يبدو مفيدًا جدًا، لكنه مع الأسف غير متسق. خذ
$!A(x) \ident \lnot x \in x$؛ عندئذ يقرر مبدأ الاستيعاب
\[
\lexists[y][\lforall[x][(x \in y \liff x \notin x)]],
\]
أي إنه يقرر وجود مجموعة جميع المجموعات غير المنتمية إلى نفسها. ولا
يمكن أن توجد مجموعة كهذه---وهذه هي مفارقة راسل. والواقع أن
$\Log{ZFC}$ تحتوي نسخة مقيدة---ومتسقة---من هذا المبدأ، هي مبدأ الفصل:
\[
\lforall[z][\lexists[y][\lforall[x][(x \in y \liff (x \in z \land
      !A(x)))]]].
\]

\begin{prob}
بين أن مبدأ الاستيعاب غير متسق بإعطاء !!a{derivation} يبين أن
\[
\lexists[y][\lforall[x][(x \in y \liff x \notin x)]] \Proves \lfalse.
\]
وقد يفيد أن تبين أولًا أن $(A \lif \lnot A) \land (\lnot A \lif A)
\Proves \lfalse$.
\end{prob}
\end{document}
